\clearpage
\SourceLeaf{207}
\renewcommand{\BookRunningTitle}{Handschriftlicher Nachlass}
\thispagestyle{empty}
\pdfbookmark[0]{Handschriftlicher Nachlass}{handschriftlicher-nachlass}
\vspace*{.25\textheight}
\begin{center}
{\Large HANDSCHRIFTLICHER}\par\vspace{2.8cm}
{\Huge NACHLASS.}
\end{center}
\vfill

\clearpage
\SourceLeaf{209}
\renewcommand{\BookRunningTitle}{Nachlass. Analysis residuorum}
\pdfbookmark[0]{Solutio congruentiae X^m-1≡0. Analysis residuorum}{analysis-residuorum-pars-prior}
\begin{center}
{\large SOLUTIO CONGRUENTIAE $X^m-1\equiv0$.}\par\medskip
ANALYSIS RESIDUORUM. CAPUT SEXTUM. PARS PRIOR.
\end{center}
\BookRule{.18\linewidth}
\renewcommand{\Article}[1]{\par\medskip\Needspace{4\baselineskip}\pdfbookmark[1]{Art. #1}{analysis-residuorum-art-#1}\centerline{\textbf{#1.}}\nobreak\smallskip}
\Article{237}
In Cap. \textsc{iii} docuimus, congruentiam $x^n\equiv1$, si pro modulo accipiatur numerus primus $p$, habere $\mu$ radices, quando $\mu$ est maxima communis mensura numerorum $n$ et $p-1$, hasque radices cum radicibus congr. $x^\mu\equiv1$ penitus convenire. Quamobrem eum casum considerare sufficit, ubi $n$ est pars aliquota numeri $p-1$. Quod autem non modo congruentiae $x^n-1\equiv0$ sed cuiusvis alius solutio pro modulis quibuscunque ex solutione pro modulis, qui sunt numeri primi, possit derivari, iam passim est ostensum infraque (Cap. \textsc{viii}) fusius docebitur.

\Article{238}
Sed ne hic quidem subsistere opus est; namque eodem Capite \textsc{iii} exposuimus, congruentiae $x^n\equiv1$ solutionem a resolutione similium congruentiarum pendere $x^a\equiv1$, $x^b\equiv1$ etc., ubi $a,b$ etc. sunt numeri primi aut numerorum primorum potestates et $n$ productum ex his numeris. Si scilicet $A,B$ etc. sunt respective radices quaecunque congruentiarum $x^a\equiv1$, $x^b\equiv1$ etc., productum ex his $AB\ldots$ erit aliqua e radicibus congruentiae $x^n\equiv1$. Nostrae igitur investigationes ad solutionem congruentiae $x^n\equiv1\;(\mathrm{mod.}\,p)$ restringentur, quando $p$ est numerus primus, $n$ numerus primus aut numeri primi potestas, simulque pars aliquota numeri $p-1$.

\SourceLeaf{210}
\Article{239}
Porro ex Cap. \textsc{iii} constat, inter congruentiae $x^n\equiv1$ radices semper aliquas dari, per quarum potestates omnes ceterae exhiberi possunt. Ita si $r$ designet huiusmodi radicem (\emph{primitivam} supra diximus, quando $n=p-1$, hancque expressionem hic quamquam significatione latiori retinebimus) omnes congr. propos. radices erunt
\[
1,\ r,\ rr,\ r^3\quad\dotsfive\quad r^{n-1}
\]
Huiusmodi ergo radices omni studio sunt investigandae, quoniam his inventis ceterae sponte patebunt. Brevitatis gratia quamcunque ipsius $r$ potestatem per exponentem uncis inclusum designamus, ita ut $(0)$ denotet unitatem, $(1)$ radicem quamcunque primitivam congruentiae $x^n\equiv1$, $(2)$ ipsius $(1)$ quadratum etc.: ita ut haec series $(0),\ (1),\ (2),\ (3),\dotsfour(n-1)$ omnes radices amplectatur. Ceterum constat, $(k)$ semper fore talem radicem primitivam, quoties $k$ ad $n$ est primus; i.~e. nostro casu (ubi $n$ est numeri primi $t$ potestas $=t^\nu$), quoties $t$ ipsum $k$ non dividit. Manifesto vero signa $(1)$, $(2)$ etc. per se sunt indeterminata; sed simulac ipsi $(1)$ valor aliquis determinatus tribuitur, omnia cetera determinata fient.

\Article{240}
Quoniam radices primitivas prae ceteris investigare propositum est, has a ceteris primum separare oportet. Quod fiet, si e serie $(0),(1),(2)\ldots(n-1)$ omnes terminos $(k)$ eiiciamus, ubi $k$ per $t$ dividitur; quodsi autem $n$ est numerus primus seu $\nu=1$, unicus $(0)$ erit abrogandus. Priusquam vero ad disquisitionem radicum superstitum progrediamur, lectorem sedulo admonemus exempla aliquot sibi conficere, ut omnia, quae sine his forsan generalius dicta viderentur, in concreto intueri possit. Nos aliquod apponimus; sed non ideo superfluum erit alia proprio Marte elaborare.

Sit $p=29$; $n=7$ et septenae congruentiae $x^7\equiv1\;(\mathrm{mod.}\,29)$ radices erunt $1,7,16,20,23,24,25$. Quoniam $n$ est numerus primus, omnes hae radices praeter $1$ erunt primitivae; posito igitur $7=(1)$ signa haec significabunt:
\[
\begin{array}{ccccccc}
(0)&(1)&(2)&(3)&(4)&(5)&(6)\\
1&7&20&24&23&16&25
\end{array}
\]

\SourceLeaf{211}
Quivis ceterum memor erit, signa $(n)$ et $(0)$, $(n+1)$ et $(1)$ etc. et in genere $(a)$ et $(b)$ aequivalere, quoties $a\equiv b\;(\mathrm{mod.}\,n)$.

\Article{241}
Sed ad nostrum propositum alio adhuc modo erit procedendum. Videlicet eos tantum terminos $(k)$ retinemus, ubi $k$ per $t$ non dividitur, quorum multitudo est $\frac{t-1}{t}\pdot n=\lambda$; omnes autem hi numeri (aut ipsis secundum $n$ congrui) per potestates successivas alicuius numeri exhiberi possunt. Sit hic $=\rho$; quare omnes radices primitivae congruentiae $x^n\equiv1$ ita denotabuntur
\[
(1)\quad(\rho)\quad(\rho^2)\quad(\rho^3)\ \dotsfour\ (\rho^{\lambda-1})
\]
Hoc autem artificio id obtinemus, ut omnes radices non primitivae penitus excludantur, cuius rei rationes et emolumenta infra clarius cognoscentur. In nostro igitur exemplo ponere possumus $\rho=3$ et radices congruentiae $x^7\equiv1$ primitivae ita ordinantur
\[
\begin{array}{rcccccc}
&(1)&(3)&(3^2)&(3^3)&(3^4)&(3^5)\\
\text{seu}&(1)&(3)&(2)&(6)&(4)&(5)\\
\text{quae erunt}&7&24&20&25&23&16
\end{array}
\]

\Article{242}
Ne lector ignarus sit, quorsum disquisitiones sequentes tendant, theorema, quod demonstrandum atque dilucidandum nobis proponimus, indicare iuvabit.

\emph{Si numerus $\lambda$ (qui est $=t^{\nu-1}\pdot t-1$) habeat factores simplices $a,b,c,d$ etc. et sit $\lambda=a^\alpha b^\beta c^\gamma\ldots$, resolutio congruentiae $x^n-1\equiv0$ pendet a resolutione $\alpha+\beta+\dotstwo$ congruentiarum inferiorum, quarum $\alpha$ sunt gradus $a$, $\beta$ gradus $b$, $\gamma$ gradus $c$ etc.}

Ita in nostro exemplo congruentiae $x^7\equiv1$ resolutio pendet a congruentia secundi gradus et ab alia tertii gradus; perspiciturque in genere numquam gradum harum congruentiarum a modulo $p$ pendere. Ut autem ad huius theorematis demonstrationem perveniamus, necesse est aliquas propositiones ad nexum inter congruentias earumque radices spectantes praemittere, quamquam proprie in Cap. octavo hae disquisitiones ulterius sint persequendae.

\SourceLeaf{212}
\Article{243}
\textsc{Theorema.}\quad\emph{Si congruentia}
\[
x^m+Ax^{m-1}+Bx^{m-2}+\dotsfour+N\equiv0\;(\mathit{mod.\ primus})
\]
\emph{ita sit comparata, ut confecto producto ex $m$ factoribus $x-r$, $x-r'$, $x-r''$, $x-r'''\ldots$ quod sit $x^m+ax^{m-1}+bx^{m-2}\ldots+n$, sit $A\equiv a$, $B\equiv b$, $C\equiv c$ etc. secundum mod.\,$p$, quantitates $r,r',r''\ldots$ erunt radices congruentiae propositae nullasque alias habebit.}

\emph{Demonstratio.}\quad I. Erit semper
\[
\begin{aligned}
&x^m+Ax^{m-1}+Bx^{m-2}+\ldots\\
&\qquad\equiv x^m+ax^{m-1}+bx^{m-2}+\ldots\;(\mathrm{mod.}\,p)
\end{aligned}
\]
Sed posterior congruentiae pars fit $=0$ ponendo $x=r$, $x=r'$, $x=r''$ etc., quare pro his ipsius $x$ valoribus prior pars fiet $\equiv0\;(\mathrm{mod.}\,p)$. Q. E. Primum.

II. Si autem alius adhuc valor $\rho$ nulli horum $r,r'$ etc. congruus congruentiae propositae satisfaceret, foret
\[
\begin{aligned}
0&\equiv\rho^m+A\rho^{m-1}+B\rho^{m-2}+\ldots\\
&\equiv\rho^m+a\rho^{m-1}+b\rho^{m-2}+\ldots\\
&\equiv(\rho-r)(\rho-r')(\rho-r'')(\rho-r''')\ldots
\end{aligned}
\]
sed quoniam nullus factorum $\rho-r$, $\rho-r'$, $\rho-r''$, etc. est $\equiv0$, productum ex omnibus fieri $\equiv0$, ob $p$ primum est absurdum. Quare praeter radices $r,r'$ etc. nullae dantur aliae. Q. E. Secundum.

\Article{244}
\textsc{Problema.}\quad\emph{Sint $r,r',r''\ldots$ quantitates incognitae, quarum multitudo sit $=m$, quarum summa sit $=\alpha$, summa quadratorum $=\beta$, summa cuborum $=\gamma\ldots$, summa potestatum, quarum exponens est $m$, $=\mu$, danturque non hi numeri (quorum multitudo etiam $=m$) ipsi, sed alii $\alpha',\beta',\gamma'$ etc. singulis congrui secundum modulum $p$, qui sit numerus primus et $>m$, invenire congruentiam $m^{\mathrm{ti}}$ gradus, cuius radices sint $r,r',r''$ etc.}

\emph{Solutio.}\quad Considerentur $r,r',r''$ etc. quasi radices alicuius aequationis
\[
x^m+Ax^{m-1}+Bx^{m-2}+Cx^{m-3}+\ldots=0
\]
determinenturque eius coëfficientes $A,B,C$ etc. (adhibendo tantummodo congruentiam loco aequalitatis) ad methodum cognitam, faciendo scilicet
\SourceLeaf{213}
\[
\begin{aligned}
-A&\equiv\alpha'\\
-2B&\equiv\beta'+A\alpha'\\
-3C&\equiv\gamma'+A\beta'+B\alpha'\\
-4D&\equiv\delta'+A\gamma'+B\beta'+C\alpha'\\
&\text{etc.}\\
-mN&\equiv\mu'+A\lambda'+\text{etc.}
\end{aligned}
\]
Hi vero coëfficientes non possunt esse indeterminati, quia omnes numeri $1,2,3\ldots m<p$. Dico congruentiam
\[
x^m+Ax^{m-1}+Bx^{m-2}+\ldots+N\equiv0
\]
esse quaesitam.

\emph{Demonstr.}\quad Ponatur aequationem, cuius radices sunt $r,r',r'',r'''$ etc., esse hanc
\[
x^m+ax^{m-1}+bx^{m-2}+\ldots=0
\]
eritque
\[
\begin{aligned}
-a&=\alpha\\
-2b&=\beta+a\alpha\\
-3c&=\gamma+a\beta+b\alpha\\
-4d&=\delta+a\gamma+b\beta+c\alpha\\
&\text{etc.}
\end{aligned}
\]
Cuique autem manifestum hinc erit, fore
\[
a\equiv A,\quad b\equiv B,\quad c\equiv C\ \text{etc.}\;(\mathrm{mod.}\,p)
\]
quare per \S\ praec. numeri $r,r',r''$ etc., qui sunt radices \emph{aequationis}
\[
x^m+ax^{m-1}+bx^{m-2}+\ldots=0
\]
erunt simul radices \emph{congruentiae}
\[
x^m+Ax^{m-1}+Bx^{m-2}+\ldots\equiv0\qquad\text{Q. E. D.}
\]
Exempla componenda lectoribus linquimus.

\Article{245}
Ad propositum nostrum revertimur. Retentis characteribus \S\S.\ 242 et antec. adhibitis ostendere aggredimur, si $\lambda$ sit productum e factoribus quibuscunque
\SourceLeaf{214}
$efg$ etc., radices congruentiae $x^n\equiv1$ primitivas, quarum multitudo est $\lambda$, ita in $e$ classes discerpi posse, ut aggregata radicum in eandem classem relatarum per congruentiam gradus $e^{\mathrm{ti}}$ dentur; his vero tamquam cognitis suppositis quamvis classem ita in $f$ ordines subdividi posse, ut aggregata cuiusvis ordinis per congruentiam $f^{\mathrm{ti}}$ gradus dentur, hique ordines rursus subdividi possunt etc., usque dum ad singulas radices perveniatur.

\Article{246}
\emph{Definitio.}\quad Complexum terminorum \emph{omnium} in tali forma $(\rho^{ke+\alpha})$ (\S.\ 241) contentorum \emph{periodum completam} sive simpliciter \emph{periodum} dicemus. Designat vero $e$ divisorem aliquem numeri $\lambda$; $\alpha$ numerum quemcunque datum, $k$ omnes numeros integros a $0$ usque ad $\frac{\lambda}{e}-1$; brevitatis vero gratia talem periodum ita designamus $(e\ast\alpha)$. Ita in exemplo nostro termini
\[
\begin{array}{rclcl}
&(1),(2),(4)&\text{periodos}&(2\ast0)&\text{constituent,}\\
&(3),(6),(5)&&(2\ast1)&\\
\text{hi vero}&(1),(6)&\text{hasce}&(3\ast0)&\\
&(3),(4)&&(3\ast1)&\\
&(2),(5)&&(3\ast2)&
\end{array}
\]

\emph{Iam si omnes termini in periodos quomodocunque distribuantur, singulaeque periodi iterum in periodos minores et sic porro, dicimus, id obtineri quod in \S.\ praec. promisimus.}

Antequam vero hanc expositionem ipsam aggrediamur, ostendemus, formationi talis periodi, quamquam a duabus quantitatibus quodammodo arbitrariis $r,\rho$ dependeat, nihil tamen vagi inesse, seu quomodocunque hae quantitates eligantur, semper eosdem terminos in eandem periodum concurrere (siquidem quot terminos periodus continere debeat, fuerit praescriptum).

Criterium, duos terminos $A,B$ in eadem periodo esse, inde petitur, quod uterque in tali forma continetur: $(\rho^{ke+\alpha})$ sive esse $A\equiv r^{\rho^{ke+\alpha}}$, $B\equiv r^{\rho^{k'e+\alpha}}\;(\mathrm{mod.}\,p)$. Hic autem $r$ est radix primitiva congruentiae $x^n\equiv1\;(\mathrm{mod.}\,p)$; $\rho$ vero radix primitiva congruentiae $x^\lambda\equiv1\;(\mathrm{mod.}\,n)$; vide supra.

Demonstrandum est, si loco numerorum $r,\rho$ alii eligantur, puta $s,\sigma$, tunc $A$ et $B$ in similibus formis $s^{\sigma^{le+\beta}}$, $s^{\sigma^{l'e+\beta}}$ comprehendi.

\SourceLeaf{215}
Sit $s^m\equiv r\;(\mathrm{mod.}\,p)$; $\sigma^\mu\equiv\rho\;(\mathrm{mod.}\,n)$ et $m\equiv\sigma^\zeta\;(\mathrm{mod.}\,n)$, quod fieri potest, quia $r,\rho$ sunt radices primitivae: erit vero $m$ primus ad $n$, $\mu$ ad $\lambda$ (Cap. \textsc{iii}). Per debitas substitutiones obtinebimus
\[
A\equiv s^{\sigma^{\mu ke+\mu\alpha+\zeta}},\qquad B=s^{\sigma^{\mu k'e+\mu\alpha+\zeta}}\qquad\text{Q. E. D.}
\]
